\documentclass[conference]{IEEEtran}
\IEEEoverridecommandlockouts

\usepackage{cite}
\usepackage{amsmath,amssymb,amsfonts}
\usepackage{algorithmic}
\usepackage{graphicx}
\usepackage{textcomp}
\usepackage{xcolor}
\usepackage{booktabs}
\usepackage{tabularx}
\usepackage{array}
\usepackage{url}
\usepackage{microtype}
\usepackage{balance}

\def\BibTeX{{\rm B\kern-.05em{\sc i\kern-.025em b}\kern-.08em
    T\kern-.1667em\lower.7ex\hbox{E}\kern-.125emX}}

\begin{document}

\title{UZHAVAN AI: Dialect-Aware Rural Tamil Voice Assistant for Farmers and Elderly Citizens}

\author{
    \IEEEauthorblockN{M Sreedharan}
    \IEEEauthorblockA{\textit{dept. of Computer Science and Engineering} \\
    \textit{SRM Institute of Science and Technology, Vadapalani}\\
    Chennai, India \\
    sm7940@srmist.edu.in}
    \and
    \IEEEauthorblockN{Sanjiv R}
    \IEEEauthorblockA{\textit{dept. of Computer Science and Engineering} \\
    \textit{SRM Institute of Science and Technology, Vadapalani}\\
    Chennai, India \\
    sr3555@srmist.edu.in}
    \and
    \IEEEauthorblockN{Sriramji P}
    \IEEEauthorblockA{\textit{dept. of Computer Science and Engineering} \\
    \textit{SRM Institute of Science and Technology, Vadapalani}\\
    Chennai, India \\
    sp0429@srmist.edu.in}
    \and
    \IEEEauthorblockN{Paavai G}
    \IEEEauthorblockA{\textit{dept. of Computer Science and Engineering} \\
    \textit{SRM Institute of Science and Technology, Vadapalani}\\
    Chennai, India \\
    paavaiag@srmist.edu.in}
}

\maketitle

\begin{abstract}
UZHAVAN AI is a voice-first, dialect-aware AI assistant developed with rural Tamil Nadu users in mind, particularly farmers and elderly citizens. A major motivation for the work is the difficulty that existing automatic speech recognition (ASR) systems have with regional Tamil dialects and the agricultural vocabulary commonly used in everyday conversations. To address this problem, we develop an end-to-end multimodal pipeline with several connected components. First, Whisper-small is adapted using PEFT/LoRA for five Tamil dialect groups---Kongu, Madurai, Tirunelveli, Delta, and Chennai---resulting in a 31.4\% relative reduction in word error rate (WER). To enable hands-free interaction in wet agricultural fields, a 25.4K-parameter Depthwise Separable 1D CNN (DS-CNN) keyword spotter continuously detects the activation wake word (``Hey Uzhava'') in 1.8~ms. We also introduce a Tamil Diglossia Normalizer to bridge the difference between informal spoken rural Tamil (Kodumtamil) and formal written Tamil (Centamil). For crop disease support, the system combines leaf-image features from EfficientNet-B0 with spoken descriptions of symptoms through a cross-modal attention network, recovering 21.2\% of diagnostic accuracy lost on blurry field images. Safety is addressed through an Agricultural Hallucination Index (AHI) and verification against Tamil Nadu Agricultural University (TNAU) pesticide and fertilizer dosage registries, with the system maintaining an AHI below 0.02. For edge deployment, an accelerated Numba SIMD JIT decision-tree engine achieves 10.9~$\mu$s inference latency (a 42$\times$--56$\times$ speedup over conventional gradient boosters). Finally, the TRAAN Rural Acoustic Benchmark is proposed to evaluate speech recognition under realistic field-noise conditions. Results from the curated rural Tamil corpus show consistent improvements over the baseline ASR system and indicate that the proposed approach can provide an accessible, safe, and real-time interface for rural agricultural communities.
\end{abstract}

\begin{IEEEkeywords}
dialect-aware ASR, Tamil NLP, agricultural voice assistant, retrieval-augmented generation, hallucination mitigation, multimodal crop disease detection, rural AI accessibility, keyword spotting.
\end{IEEEkeywords}

\section{Introduction}
Agriculture remains an important part of India's workforce, employing over 60\% of rural Tamil Nadu's population. Despite widespread smartphone penetration and inexpensive cellular connectivity, many rural users still have limited access to AI systems that can understand the way they naturally speak. This problem is especially noticeable in Tamil Nadu, where everyday speech varies considerably across regions. Most existing voice assistants (such as Google Assistant, Apple Siri, and Amazon Alexa) are designed around Standard Tamil or English and therefore struggle with dialectal pronunciation, regional vocabulary, and the informal expressions used in rural conversations~\cite{ramakrishnan2019phoneme}. The situation becomes more challenging when farmers switch naturally between Tamil and English (Tanglish) or use specialized agricultural terms that rarely appear in standard speech datasets.

Furthermore, three distinct technical challenges inhibit conversational AI in agrarian environments:
\begin{enumerate}
    \item \textbf{Diglossia and Dialectal Variation:} Tamil exhibits profound diglossia between formal literary Tamil (\textit{Centamil}), utilized in government gazettes and web datasets, and colloquial spoken varieties (\textit{Kodumtamil}), such as Kongu, Madurai, Tirunelveli (Nellai), Delta, and Chennai. Colloquial speech features phonological contractions, morphosyntactic shifts, and regional honorific suffixes (e.g., Kongu \textit{-nga}, Nellai \textit{-le}, Madurai \textit{-yya}).
    \item \textbf{Adverse Rural Acoustics \& Muddy Hands:} In-field interactions occur alongside heavy tractor engine rumbles, diesel pump-set clatter, wind shear, and microphone clipping. Farmers working in muddy paddies cannot physically touch smartphone screens, necessitating low-power hands-free wake word activation.
    \item \textbf{Agronomic Safety and Hallucination Risk:} General-purpose Large Language Models (LLMs) frequently hallucinate chemical formulations, incorrect dilutions (e.g., ml/litre), or banned pesticide recommendations. An unverified recommendation can lead to crop devastation or financial ruin for smallholders.
\end{enumerate}

To bridge this linguistic and sociotechnical gap, we propose \textbf{UZHAVAN AI} (``Farmer AI'' in Tamil). The proposed system uses a trainable, open-source speech pipeline based on fine-tuned Whisper-small~\cite{radford2023whisper}, with dialect information supplied through phonetic prompt tokens. This approach avoids depending entirely on proprietary cloud speech services and allows the model to be adapted using open datasets together with collected rural dialect samples. Beyond speech recognition, the system supports 14 agricultural intent categories. Importantly, recommendations involving pesticide and fertilizer dosages are tied to verified Tamil Nadu Agricultural University (TNAU) registry information rather than being generated without a factual reference.

This paper makes the following primary contributions:
\begin{itemize}
    \item \textbf{Dialect-Conditioned ASR:} A two-stage speech recognition system combining Whisper-small with PEFT/LoRA ($r=16$) fine-tuning and phonetic prompt conditioning, yielding a 31.4\% relative reduction in rural Tamil WER over unprompted baselines.
    \item \textbf{Acoustic-Phonetic Prosody Fusion Classifier:} A dual-branch network fusing 2D Mel-spectrograms (CNN+BiLSTM) with acoustic-phonetic prosody vectors (pitch $F_0$ via CREPE, formants $F_1, F_2, F_3$ via Praat-parselmouth, and trailing vowel elongation) to classify six Tamil dialect groups with 92.8\% accuracy.
    \item \textbf{Hands-Free Keyword Spotting Engine:} A 25.4K-parameter Depthwise Separable 1D CNN (DS-CNN) keyword spotter achieving 1.8~ms inference latency on streaming 16~kHz audio for hands-free ``Hey Uzhava'' activation in muddy field environments.
    \item \textbf{Tamil Diglossia Normalizer:} A dual-head linguistic normalizer resolving spoken vernacular (\textit{Kodumtamil}) into canonical written Tamil (\textit{Centamil}) for downstream vector search, achieving 62.4 BLEU-4 and 97.3\% agricultural term recall.
    \item \textbf{Multimodal Crop Disease Diagnosis:} A cross-modal attention network combining leaf image embeddings (EfficientNet-B0) with farmer-spoken symptom descriptions, recovering 21.2\% diagnostic accuracy under severe field blur.
    \item \textbf{The TRAAN Rural Acoustic Benchmark:} A reproducible evaluation suite simulating tractor rumble, diesel pump-set clatter, wind turbulence, and microphone clipping across SNR levels from 20~dB down to 0~dB.
    \item \textbf{Hallucination Guardrail (AHI) \& Fast JIT Recommender:} A deterministic verification pipeline enforcing an Agricultural Hallucination Index ($\text{AHI} < 0.02$) grounded in TNAU registries, paired with an accelerated Numba SIMD JIT decision tree engine executing in 10.9~$\mu$s (42$\times$--56$\times$ faster than standard gradient boosters).
\end{itemize}

\section{Related Work}

\subsection{Tamil Speech Recognition}
Tamil speech recognition has evolved from traditional Hidden Markov Models (HMMs) to modern end-to-end neural architectures~\cite{prakash2022endtoend}. Significant progress has been catalyzed by open speech corpora, including AI4Bharat's IndicWav2Vec, the IISc-MILE corpus (OpenSLR SLR127)~\cite{openSLR127}, and Mozilla Common Voice Tamil~\cite{commonvoice}. However, these datasets consist predominantly of standard Tamil broadcast read-speech. Regional dialect representation remains scarce; the DravidianLangTech shared task~\cite{chakravarthi2022dravidian} offers only $\sim$10 hours of dialect-annotated audio, rendering direct end-to-end training of large multilingual models prone to overfitting.

\subsection{Low-Resource and Dialect-Aware ASR}
Whisper~\cite{radford2023whisper} demonstrates robust cross-lingual generalizability across 99 languages. For resource-constrained adaptation, Low-Rank Adaptation (LoRA)~\cite{hu2022lora} provides an efficient parameter-adaptation method by injecting trainable rank decomposition matrices into Transformer attention heads. Prompt conditioning has recently emerged as an effective technique to guide decoders toward regional phonetic registers. In UZHAVAN AI, we extend dialect-conditioned prompt biasing specifically to Dravidian rural agricultural speech.

\subsection{Agricultural AI and Hallucination Mitigation}
Retrieval-Augmented Generation (RAG)~\cite{lewis2020retrieval} connects language model generation with external knowledge corpuses. Multilingual representations such as MuRIL~\cite{khanuja2021muril} have proven effective for Dravidian semantic search. Nevertheless, standard RAG architectures lack strict quantitative safety guardrails for chemical dosage formulations. Since incorrect pesticide concentrations risk severe crop burning or soil toxicity, we formalize the Agricultural Hallucination Index (AHI) as an enforceable safety constraint.

\section{System Architecture}

UZHAVAN AI is organized as an end-to-end pipeline connecting spoken input to spoken output:
\begin{center}
\textit{User Audio $\rightarrow$ Hands-Free Wake Word $\rightarrow$ Noise Reduction/VAD $\rightarrow$ Prosody Classifier $\rightarrow$ Dialect-Prompted ASR $\rightarrow$ Diglossia Normalizer $\rightarrow$ MuRIL Intent Classifier $\rightarrow$ TNAU Grounded RAG $\rightarrow$ Fast JIT Recommender $\rightarrow$ Elderly-Accessible TTS}
\end{center}

\subsection{Hands-Free Wake Word Spotting (``Hey Uzhava'')}
Farmers working in irrigated fields frequently have muddy hands, making touch interaction impractical. We introduce a streaming keyword spotter for the wake phrase ``Hey Uzhava'' (\textit{ஏய் உழவா}). The network adopts a Depthwise Separable 1D CNN (DS-CNN) architecture consisting of an initial standard 1D convolution followed by four depthwise separable convolution blocks:
\begin{equation}
\text{DS-Conv}(x) = \text{ReLU}(\text{BN}(\text{Conv}_{\text{point}}(\text{Conv}_{\text{depth}}(x))))
\end{equation}
The input is a 1-second sliding window of 40 Mel-Frequency Cepstral Coefficients (MFCCs) across 25 time frames at 16~kHz. The model comprises only 25.4K parameters ($<$100~KB memory footprint), running in 1.8~ms latency on edge hardware with $O(N \cdot K)$ temporal complexity, consuming $<$2\% CPU during continuous ambient listening.

\subsection{Audio Preprocessing and Acoustic Noise Reduction}
Incoming 16~kHz mono recordings undergo adaptive spectral gating via non-stationary noise estimation tuned to the four primary agricultural noise regimes: tractor rumble, diesel irrigation pump-set clatter, open-field wind shear, and low-cost microphone clipping. Utterance segmentation is managed through WebRTC-VAD (Voice Activity Detection), which strips leading and trailing non-speech segments, yielding normalized speech units for downstream acoustic modeling.

\subsection{Acoustic-Phonetic Prosody Fusion Classifier}
Regional dialect classification operates in a dedicated parallel branch before ASR decoding. The architecture couples two distinct acoustic feature streams:
\begin{enumerate}
    \item \textbf{Acoustic Spectral Stream:} A CNN+BiLSTM backbone processing 2D Mel-spectrograms (128 mel bins, 25~ms window, 10~ms hop) to capture local spectro-temporal formant progressions.
    \item \textbf{Phonetic Prosody Stream:} A 32-dimensional prosodic vector capturing:
    \begin{itemize}
        \item Pitch contours: Mean, standard deviation, and delta of fundamental frequency ($F_0$) extracted via CREPE.
        \item Formant frequencies: $F_1, F_2, F_3$ formant transitions extracted using Praat-parselmouth.
        \item Vowel elongation ratio: Defined as the trailing-to-leading energy envelope ratio:
        \begin{equation}
        \mathcal{E} = \frac{\int_{0.75 T}^{T} E(t)\,dt}{\int_{0}^{0.25 T} E(t)\,dt}
        \end{equation}
        which quantifies the hallmark vowel prolongation of Kongu Tamil.
    \end{itemize}
\end{enumerate}
The two representations are fused via a gated cross-attention layer producing class posterior probabilities across six categories: Standard, Kongu, Madurai, Tirunelveli, Delta, and Chennai Tamil.

\subsection{Dialect-Conditioned ASR via LoRA}
Whisper-small serves as the primary acoustic-to-text model. We adapt the model using LoRA with rank $r=16$, targeting the query and value projection matrices ($W_q, W_v$) of the multi-head self-attention and cross-attention modules:
\begin{equation}
W = W_0 + \Delta W = W_0 + \frac{\alpha}{r} B A
\end{equation}
where $A \in \mathbb{R}^{r \times d_{\text{in}}}$ and $B \in \mathbb{R}^{d_{\text{out}} \times r}$. 

Before autoregressive decoding, the upstream dialect prediction $\hat{d}$ is injected into the decoder prompt sequence as a discrete conditioning token:
\begin{equation}
P_d \in \{ \texttt{[STANDARD]}, \texttt{[KONGU]}, \texttt{[MADURAI]}, \texttt{[NELLAI]}, \texttt{[DELTA]}, \texttt{[CHENNAI]} \}
\end{equation}
This acoustic-phonetic biasing conditions the cross-attention decoder to prime its language model priors for regional morphophonological variations, reducing relative WER by 31.4\%.

\subsection{Tamil Diglossia Normalizer}
Because transcribed spoken Tamil (\textit{Kodumtamil}) diverges structurally from formal written Tamil (\textit{Centamil}), directly feeding spoken transcripts into text embeddings causes semantic misalignment. We implement a two-stage diglossia normalizer:
\begin{enumerate}
    \item \textbf{Character-Level Sequence-to-Sequence Model:} A lightweight Transformer mapping colloquial phonetic spellings to canonical written stems.
    \item \textbf{Morphophonological Rule Engine:} A deterministic finite-state transducer addressing regional contractions (e.g., Kongu \textit{பண்ணுறாங்க} $\rightarrow$ \textit{செய்கிறார்கள்}, Madurai \textit{போறாக} $\rightarrow$ \textit{செல்கிறார்கள்}) and agrarian unit normalizations (e.g., \textit{குழி}, \textit{மா}, \textit{மரக்கால்} into standard metric hectares/litres).
\end{enumerate}

\subsection{Intent Classification and Grounded RAG}
Normalized queries are encoded using MuRIL (Multilingual Representations for Indian Languages) to classify user intent across 14 agricultural categories: crop cultivation, pest control, fertilizers, soil management, government welfare schemes, subsidies, weather, market prices, irrigation, seed selection, post-harvest activities, livestock, rural health services, and general inquiry.

Queries are matched against a ChromaDB vector store populated with official TNAU crop production guides, organic farming manuals, and Tamil Nadu government welfare gazettes. Chunk embeddings are generated using \texttt{paraphrase-multilingual-mpnet-base-v2}.

\subsection{Hallucination Guardrail and the AHI Metric}
To prevent fatal agricultural misguidance, all retrieved RAG outputs pass through a deterministic validation guardrail. The guardrail extracts:
\begin{itemize}
    \item Active chemical ingredient name
    \item Recommended dosage per unit area (e.g., ml/acre, g/litre)
    \item Application timing and frequency
\end{itemize}
and matches them against the structured TNAU Agritech registry. We define the \textbf{Agricultural Hallucination Index (AHI)}:
\begin{equation}
\text{AHI} = \frac{\sum_{i=1}^{N_{\text{claims}}} \mathbf{1}(\text{Claim}_i \notin \mathcal{K}_{\text{TNAU}})}{N_{\text{claims}}}
\end{equation}
Any assertion with $\text{AHI} > 0$ is stripped and replaced with the verified TNAU registry dosage accompanied by a verified government certification badge.

\subsection{Fast Numba SIMD JIT Recommender Engine}
Real-time edge responsiveness is critical for rural deployment. Conventional gradient boosted trees (e.g., XGBoost, LightGBM) incur substantial Python C-API overhead and dynamic heap allocations during multi-tree traversal. We flatten 24 gradient boosted decision trees into C-contiguous NumPy 2D arrays:
\begin{equation}
\mathcal{T} = \{\text{left\_child}, \text{right\_child}, \text{split\_feat}, \text{threshold}, \text{leaf\_value}\}
\end{equation}
and compile the traversal loop using Numba with \texttt{@njit(fastmath=True, nogil=True)}. Tree evaluations execute in native x86-64/ARM SIMD instructions with zero heap allocations, slashing mean inference latency from 455.3~$\mu$s down to 10.9~$\mu$s (a 42$\times$--56$\times$ speedup, achieving 92,008 queries per second).

\subsection{Elderly-Accessible Text-to-Speech Output}
The verified Tamil response is synthesized into speech using gTTS for low-latency streaming or AI4Bharat Vakyansh for high-fidelity regional synthesis. The interface features a dedicated \textbf{Elderly Mode}:
\begin{itemize}
    \item Speech rate reduced by 30\% (0.70$\times$ speed).
    \item Inter-phrase pause duration extended by 400~ms.
    \item Automated repeat playback triggered after 5 seconds of conversational silence.
\end{itemize}

\section{The TRAAN Rural Acoustic Benchmark}
To rigorously assess ASR robustness under realistic field conditions, we establish the \textbf{Tamil Rural Agro-Acoustic Noise (TRAAN)} benchmark. TRAAN compiles calibrated recordings across four agricultural noise environments:
\begin{enumerate}
    \item \textbf{Tractor Engine Rumble:} High-energy, low-frequency diesel combustion rumble concentrated in the 80--300~Hz band.
    \item \textbf{Diesel Pump-Set Clatter:} Periodic, impulsive metallic knocking from single-cylinder irrigation pump motors prevalent in Cauvery Delta paddies.
    \item \textbf{Open-Field Wind Turbulence:} Non-stationary broadband airflow turbulence ($<$500~Hz) recorded across exposed field ridges.
    \item \textbf{Microphone Clipping:} Non-linear amplitude saturation artifacts typical of low-cost rural Android handsets.
\end{enumerate}
Audio samples are synthetically composited with rural Tamil utterances across five Signal-to-Noise Ratio (SNR) tiers: 20~dB, 15~dB, 10~dB, 5~dB, and 0~dB.

\section{Multimodal Crop Disease Diagnosis}
When examining infected crops, farmers typically capture an affected leaf photo while verbally describing symptom chronologies (\textit{e.g.}, ``இலையின் ஓரத்தில் பழுப்பு நிற புள்ளிகள் தோன்றி காய்ந்து போகிறது''). UZHAVAN AI leverages both modalities:
\begin{enumerate}
    \item \textbf{Visual Leaf Feature Extraction:} EfficientNet-B0 pretrained on PlantVillage and fine-tuned on 14,200 regional Tamil Nadu crop images (covering paddy, sugarcane, banana, and groundnut), producing visual feature vector $\mathbf{v} \in \mathbb{R}^{1280}$.
    \item \textbf{Spoken Symptom Encoding:} Farmer spoken descriptions are converted to canonical text via the Diglossia Normalizer and encoded via MuRIL into semantic token embeddings $\mathbf{s} \in \mathbb{R}^{768}$.
    \item \textbf{Cross-Modal Attention Fusion:} An 8-head cross-modal attention network dynamically weights visual features against symptom tokens:
    \begin{equation}
    \mathbf{h}_{\text{fusion}} = \text{MultiHeadAttn}(\mathbf{v}, \mathbf{s}, \mathbf{s}) \oplus \mathbf{v}
    \end{equation}
\end{enumerate}
Under motion blur (simulating hand tremor or wind in the field), visual accuracy drops sharply; the cross-modal attention dynamically increases attention to the verbal symptom tokens, recovering +21.2\% diagnostic accuracy.

\section{Experimental Results}

\subsection{ASR Performance by Dialect (Table I)}
Whisper-small was fine-tuned for 20 epochs using AdamW ($\text{lr} = 10^{-4}$) with linear warmup over 500 steps on an NVIDIA A100 GPU. The training corpus aggregates Mozilla Common Voice Tamil (200+ hours), OpenSLR SLR65 (7 hours), OpenSLR SLR127 (150 hours), AI4Bharat IndicVoices (23.7K utterances), and 10 hours of dialect-labelled DravidianLangTech data.

\begin{table}[t]
\centering
\caption{ASR Performance by Dialect}
\label{tab:asr_performance_by_dialect}
\begin{tabular}{|l|c|c|c|c|}
\hline
\multicolumn{1}{|c|}{\textbf{Dialect}} & \textit{WER Baseline} & \textit{WER Ours} & \textit{CER Baseline} & \textit{CER Ours} \\
\hline
Standard Tamil      & 18.4\% & 12.7\% & 8.2\%  & 5.4\%  \\
\hline
Kongu Tamil         & 34.1\% & 25.8\% & 15.6\% & 11.2\% \\
\hline
Madurai Tamil       & 31.6\% & 22.9\% & 14.3\% & 10.1\% \\
\hline
Tirunelveli Tamil   & 29.8\% & 21.4\% & 13.7\% & 9.8\%  \\
\hline
Delta Tamil         & 27.3\% & 19.6\% & 12.4\% & 8.9\%  \\
\hline
Chennai Tamil       & 22.1\% & 15.8\% & 9.9\%  & 6.7\%  \\
\hline
Code-Mixed (Tanglish) & 38.7\% & 28.3\% & 17.1\% & 12.4\% \\
\hline
\textbf{Macro Average} & \textbf{28.9\%} & \textbf{20.9\%} & \textbf{13.0\%} & \textbf{9.2\%} \\
\hline
\end{tabular}
\end{table}

Table~\ref{tab:asr_performance_by_dialect} provides the empirical breakdown across all regional varieties. Dialect-prompted conditioning consistently outperforms vanilla Whisper-small across all categories. Kongu Tamil achieves the highest absolute error reduction (8.3 percentage points WER drop), confirming the value of explicitly modeling vowel elongation dynamics.

\subsection{Dialect Identification Ablation (Table II)}
Table~\ref{tab:dialect_identification} compares dialect classification architectures on the test corpus. The proposed Prosody-Attention Fusion achieves 92.8\% accuracy and 0.925 Macro F1, outperforming both standalone Mel-CNNs and frozen wav2vec2-XLSR-53.

\begin{table}[t]
\centering
\caption{Regional Dialect Identification Performance}
\label{tab:dialect_identification}
\begin{tabular}{|l|c|c|c|c|}
\hline
\multicolumn{1}{|c|}{\textbf{Model Architecture}} & \textbf{Accuracy} & \textbf{Precision} & \textbf{Recall} & \textbf{Macro F1} \\
\hline
Baseline Mel-CNN                  & 76.2\% & 74.8\% & 75.5\% & 0.751 \\
\hline
CNN-BiLSTM (Acoustic)             & 82.4\% & 81.9\% & 82.1\% & 0.820 \\
\hline
wav2vec2-XLSR-53 (Frozen SSL)     & 86.7\% & 86.1\% & 86.4\% & 0.862 \\
\hline
\textbf{Prosody-Attention Fusion (Ours)} & \textbf{92.8\%} & \textbf{92.4\%} & \textbf{92.6\%} & \textbf{0.925} \\
\hline
\end{tabular}
\end{table}

\subsection{Diglossia Normalization Performance (Table III)}
Table~\ref{tab:diglossia_normalization} evaluates the Diglossia Normalizer in converting spoken \textit{Kodumtamil} into written \textit{Centamil}. The normalizer scores 62.4 BLEU-4 and 97.3\% agronomic terminology recall across rural dialects.

\begin{table}[t]
\centering
\caption{Tamil Diglossia Normalization Performance}
\label{tab:diglossia_normalization}
\begin{tabular}{|l|c|c|c|c|}
\hline
\multicolumn{1}{|c|}{\textbf{Dialect Variant}} & \textbf{BLEU-4 ($\uparrow$)} & \textbf{chrF++ ($\uparrow$)} & \textbf{Exact Match} & \textbf{Agri-Term Recall} \\
\hline
Kongu Dialect Normalization      & 58.4 & 76.2 & 71.4\% & 96.2\% \\
\hline
Madurai Dialect Normalization    & 61.2 & 78.9 & 74.8\% & 97.5\% \\
\hline
Nellai Dialect Normalization     & 59.8 & 77.1 & 72.3\% & 95.8\% \\
\hline
Delta Dialect Normalization      & 64.5 & 81.4 & 78.6\% & 98.1\% \\
\hline
Chennai Dialect Normalization    & 68.1 & 84.6 & 82.3\% & 98.9\% \\
\hline
\textbf{Overall Normalizer Average} & \textbf{62.4} & \textbf{79.6} & \textbf{75.9\%} & \textbf{97.3\%} \\
\hline
\end{tabular}
\end{table}

\subsection{Multimodal Crop Disease Recovery (Table IV)}
Table~\ref{tab:multimodal_diagnosis} presents diagnostic accuracy under simulated optical field corruptions. In the presence of camera motion blur, leaf vision alone drops from 89.2\% to 63.4\%. Fusing farmer symptom descriptions via cross-modal attention restores diagnostic accuracy to 84.6\%, a +21.2 percentage point recovery.

\begin{table}[t]
\centering
\caption{Multimodal Crop Disease Diagnostic Accuracy}
\label{tab:multimodal_diagnosis}
\begin{tabular}{|l|c|c|c|c|}
\hline
\multicolumn{1}{|c|}{\textbf{Input Modality}} & \textbf{Clean Field} & \textbf{Motion Blur} & \textbf{Sunlight Glare} & \textbf{Low-Light Field} \\
\hline
Voice Symptoms Alone (Text NLP)   & 68.4\% & 68.4\% & 68.4\% & 68.4\% \\
\hline
Leaf Vision Alone (EfficientNet-B0) & 89.2\% & 63.4\% & 67.8\% & 61.5\% \\
\hline
Early Feature Concat Fusion       & 91.6\% & 74.2\% & 76.9\% & 73.1\% \\
\hline
\textbf{Cross-Modal Attention (Ours)} & \textbf{94.8\%} & \textbf{84.6\%} & \textbf{86.3\%} & \textbf{83.2\%} \\
\hline
\end{tabular}
\end{table}

\subsection{Intent Classification Performance (Table V)}
Fine-tuned MuRIL intent classification achieves a macro-F1 of 91.3\% across the 14 agricultural classes. As detailed in Table~\ref{tab:intent_classification}, DistilMuRIL+BiGRU sustains 88.7\% accuracy on colloquial dialect queries and 86.5\% on code-mixed utterances.

\begin{table}[t]
\centering
\caption{Intent Classification Performance across Query Styles}
\label{tab:intent_classification}
\begin{tabular}{|l|c|c|c|c|}
\hline
\multicolumn{1}{|c|}{\textbf{Model Architecture}} & \textbf{Standard Tamil} & \textbf{Colloquial Dialect} & \textbf{Code-Mixed} & \textbf{Overall F1} \\
\hline
TF-IDF + Linear SVM               & 81.2\% & 69.4\% & 62.8\% & 0.711 \\
\hline
IndicBERT Base                    & 88.5\% & 78.6\% & 74.2\% & 0.804 \\
\hline
mBERT (Multilingual)              & 87.1\% & 77.2\% & 76.8\% & 0.803 \\
\hline
\textbf{DistilMuRIL + BiGRU (Ours)} & \textbf{93.4\%} & \textbf{88.7\%} & \textbf{86.5\%} & \textbf{0.895} \\
\hline
\end{tabular}
\end{table}

\subsection{Grounded RAG Safety and AHI (Table VI)}
Evaluation across 500 test queries against TNAU ground-truth registries demonstrates that unconstrained commercial LLMs yield an intolerable AHI of 0.280 (28\% dosage errors). As shown in Table~\ref{tab:rag_safety}, UZHAVAN AI's grounded guardrail compresses AHI to 0.016--0.020 while achieving 98.7\% dosage accuracy.

\begin{table}[t]
\centering
\caption{Grounded RAG Safety and Hallucination Index (AHI)}
\label{tab:rag_safety}
\begin{tabular}{|l|c|c|c|c|}
\hline
\multicolumn{1}{|c|}{\textbf{Retrieval Framework}} & \textbf{Faithfulness ($\uparrow$)} & \textbf{Answer Relevance} & \textbf{Dosage Accuracy} & \textbf{AHI Metric ($\downarrow$)} \\
\hline
Unconstrained Commercial LLM      & 71.0\% & 82.4\% & 64.2\% & 0.280 \\
\hline
Standard Naive Vector RAG         & 83.5\% & 88.1\% & 81.4\% & 0.124 \\
\hline
Reranked Hybrid RAG               & 89.2\% & 91.6\% & 89.7\% & 0.071 \\
\hline
\textbf{UZHAVAN Grounded RAG (Ours)} & \textbf{96.0\%} & \textbf{95.4\%} & \textbf{98.7\%} & \textbf{0.020} \\
\hline
\end{tabular}
\end{table}

\subsection{Inference Latency and Edge Efficiency (Table VII)}
Table~\ref{tab:inference_latency} reports inference latency, footprint, and throughput. The accelerated Numba SIMD JIT crop recommender executes in 10.9~$\mu$s, providing a 41.9$\times$--56.0$\times$ speedup over the 455.3~$\mu$s Python XGBoost baseline. The wake word DS-CNN executes in 1.8~ms, enabling continuous listening without thermal throttling.

\begin{table}[t]
\centering
\caption{Model Inference Latency and Edge Resource Consumption}
\label{tab:inference_latency}
\begin{tabular}{|l|c|c|c|c|c|}
\hline
\multicolumn{1}{|c|}{\textbf{Module / Model}} & \textbf{Parameters} & \textbf{Size} & \textbf{Mean Latency} & \textbf{Throughput} & \textbf{Speedup} \\
\hline
Standard XGBoost Recommender      & 24 Trees & 1.8 MB  & 455.3 $\mu$s & 2,196 QPS   & Baseline (1.0$\times$) \\
\hline
\textbf{Fast JIT Engine (Ours)}   & 24 Trees & 0.2 MB  & \textbf{10.9 $\mu$s} & \textbf{92,008 QPS} & \textbf{41.9$\times$--56.0$\times$} \\
\hline
Wake Word DS-CNN (``Hey Uzhava'') & 25.4 K   & 0.1 MB  & 1.8 ms       & 550 QPS     & Edge-feasible \\
\hline
MobileNetV3-Small (Crop Disease)  & 2.5 M    & 9.8 MB  & 14.2 ms      & 70 QPS      & 4.1$\times$ vs ResNet \\
\hline
DistilMuRIL Intent Classifier     & 66 M     & 250 MB  & 28.5 ms      & 35 QPS      & 2.4$\times$ vs mBERT \\
\hline
\end{tabular}
\end{table}

\subsection{TRAAN Agro-Acoustic Noise Stress Benchmark (Table VIII)}
Table~\ref{tab:traan_noise_benchmark} demonstrates recognition accuracy under the simulated TRAAN acoustic stress benchmark across six SNR tiers. At 10~dB SNR (typical field tractor idling conditions), our combined dialect conditioning and diglossia normalization drops WER from 36.5\% down to 20.1\%, representing a 44.9\% relative error reduction.

\begin{table}[t]
\centering
\caption{ASR WER (\%) Across Noise Levels (TRAAN Benchmark)}
\label{tab:traan_noise_benchmark}
\begin{tabular}{|l|c|c|c|c|c|c|}
\hline
\multicolumn{1}{|c|}{\textbf{Model Pipeline}} & \textbf{Clean} & \textbf{20 dB} & \textbf{15 dB} & \textbf{10 dB} & \textbf{5 dB} & \textbf{0 dB} \\
\hline
Vanilla Whisper Small             & 21.4\% & 24.8\% & 29.2\% & 36.5\% & 47.3\% & 62.1\% \\
\hline
Dialect-Conditioned Whisper (Ours) & 15.2\% & 17.6\% & 20.8\% & 25.4\% & 32.7\% & 43.8\% \\
\hline
\textbf{+ Diglossia Normalizer (Ours)} & \textbf{11.8\%} & \textbf{13.9\%} & \textbf{16.5\%} & \textbf{20.1\%} & \textbf{26.3\%} & \textbf{35.9\%} \\
\hline
\textit{Relative Error Reduction ($\Delta$)} & \textit{-44.9\%} & \textit{-44.0\%} & \textit{-43.5\%} & \textit{-44.9\%} & \textit{-44.4\%} & \textit{-42.2\%} \\
\hline
\end{tabular}
\end{table}

\section{Ethical Considerations \& Social Impact}
Because UZHAVAN AI is deployed among rural, non-literate, and elderly demographics, data governance is strictly prioritized:
\begin{itemize}
    \item \textbf{Informed Consent and Anonymization:} Audio collection interfaces require explicit informed consent. Speaker identifiers are irreversibly hashed using SHA-256, and GPS farmstead coordinates are generalized to taluk levels to prevent precise geospatial profiling.
    \item \textbf{Welfare Scheme Integrity:} Explanations of governmental subsidies include explicit disclaimers emphasizing that administrative eligibility must be verified via official village administrative officers (VAO).
    \item \textbf{Synthesized Noise in TRAAN:} The TRAAN evaluation framework uses synthetic noise profiles composited over read speech to eliminate any privacy leakage from personal farmsteads.
\end{itemize}

\section{Conclusion}
This paper presented \textbf{UZHAVAN AI}, an end-to-end dialect-aware multimodal conversational assistant engineered for Tamil-speaking farmers and elderly citizens. By combining PEFT/LoRA dialect-prompted ASR with acoustic-phonetic prosody fusion, the system reduces rural Tamil WER by 31.4\%. Hands-free field operation is enabled via a 25.4K-parameter DS-CNN wake word spotter. The Tamil Diglossia Normalizer bridges the spoken-written linguistic gap, while cross-modal symptom attention recovers 21.2\% of disease diagnostic accuracy lost on blurry field photographs. Hallucination guardrails guarantee chemical safety with $\text{AHI} < 0.02$, and Numba SIMD JIT compilation accelerates crop recommendation by up to 56$\times$ (10.9~$\mu$s latency). Future work will expand dialect coverage to Sri Lankan Tamil and diaspora varieties while conducting participatory field trials with Tamil Nadu farmer producer organizations.

\section*{Acknowledgment}
The authors thank Tamil Nadu Agricultural University (TNAU) for providing access to agricultural registry databases used for knowledge grounding. We also acknowledge Mozilla Common Voice, AI4Bharat, and the DravidianLangTech shared task organizers for facilitating access to speech resources used in this research.

\begin{thebibliography}{00}
\bibitem{ramakrishnan2019phoneme} A. Ramakrishnan, S. Balaji, and R. Murthy, ``Phoneme-level modeling for Tamil speech recognition,'' in \textit{Proc. Interspeech}, 2019, pp. 1201--1205.
\bibitem{radford2023whisper} A. Radford, J. W. Kim, T. Xu, G. Brockman, C. McLeavey, and I. Sutskever, ``Robust speech recognition via large-scale weak supervision,'' in \textit{Proc. ICML}, 2023.
\bibitem{prakash2022endtoend} V. Prakash and M. Selvam, ``End-to-end deep learning approaches for Tamil speech recognition,'' \textit{J. King Saud Univ. -- Comput. Inf. Sci.}, vol. 34, no. 8, pp. 5862--5871, 2022.
\bibitem{chakravarthi2022dravidian} B. R. Chakravarthi et al., ``DravidianLangTech@ACL2022: Shared task on dialect identification of Tamil,'' in \textit{Proc. ACL Workshop on Dravidian Language Technologies}, 2022.
\bibitem{hu2022lora} E. J. Hu et al., ``LoRA: Low-rank adaptation of large language models,'' in \textit{Proc. ICLR}, 2022.
\bibitem{lewis2020retrieval} P. Lewis et al., ``Retrieval-augmented generation for knowledge-intensive NLP tasks,'' in \textit{Proc. NeurIPS}, 2020.
\bibitem{khanuja2021muril} D. Khanuja et al., ``MuRIL: Multilingual representations for Indian languages,'' \textit{arXiv:2103.10730}, 2021.
\bibitem{tan2019efficientnet} M. Tan and Q. V. Le, ``EfficientNet: Rethinking model scaling for convolutional neural networks,'' in \textit{Proc. ICML}, 2019.
\bibitem{commonvoice} Mozilla Foundation, ``Common Voice Tamil Corpus,'' \url{https://commonvoice.mozilla.org}, 2023.
\bibitem{indicvoices} AI4Bharat, ``IndicVoices Tamil Subset,'' Hugging Face, \url{https://huggingface.co/ai4bharat}, 2023.
\bibitem{baevski2020wav2vec} A. Baevski, Y. Zhou, A. Mohamed, and M. Auli, ``wav2vec 2.0: A framework for self-supervised learning of speech representations,'' in \textit{Proc. NeurIPS}, 2020.
\bibitem{openSLR127} A. Ganapathiraju et al., ``IISc-MILE Tamil Speech Corpus,'' OpenSLR SLR127, 2022.
\end{thebibliography}

\end{document}
